% TOP_PROBLEM: {"id":30007657,"problem_number":"LOCAL-30007657","title":"Aid that builds lasting capacity","category_id":21,"category":{"id":21,"name":"economics","display_name":"Economics","description":"Open economic theory statements, published research agendas, and proposed scoped specifications; question provenance and historical priority are qualified per record.","slug":"economics","order_index":21,"created_at":"2026-10-08T00:00:00Z"},"status":"provenance_pending","external_url":null,"legacy_ids":[],"published":false,"statement":"In the explicitly specified setting: Which aid designs produce durable economic gains and stronger local institutions after funding ends, and which create dependence or substitute for domestic capacity? The mathematical statement above fixes the target, assumptions and scope of this catalogue formulation.","economics":{"original_id":146,"field":"Development and poverty","kind":"identification","formalization_basis":"proposed_specification","source_status":"not_verified","priority_status":"not_established","priority_note":"No qualifying problem statement was directly verified, so historical priority cannot be assigned.","earliest_verified_reference":null,"current_reference":{"authors":["Katherine Casey","Rachel Glennerster","Edward Miguel","Maarten J. Voors"],"title":"Long Run Effects of Aid: Forecasts and Evidence from Sierra Leone","year":2021,"url":"https://www.nber.org/papers/w29079","doi":"10.3386/w29079","locator":"Abstract, paragraph describing eleven-year follow-up and institutional outcomes","evidence_summary":"This follow-up supplies persistent infrastructure and market results but limited institutional change; the broad aid-versus-state-capacity question is not explicitly declared open."},"formalization_note":"This is a proposed scoped research specification. No primary passage explicitly stating this entire remaining question as open was verified. The supporting literature may answer important subquestions; this entry must not be advertised as a verified formal open problem."},"statusQualification":"Provenance pending: an explicit published statement of this exact open question was not verified. This is a catalogue-authored candidate specification, not a certified open conjecture. This is a proposed scoped research specification. No primary passage explicitly stating this entire remaining question as open was verified. The supporting literature may answer important subquestions; this entry must not be advertised as a verified formal open problem.","statusReviewedAt":"2026-10-08"}
% FIRST_POSTED: 2026-10-08
% ENTRY_KIND: statement_only
% !TeX program = lualatex
\documentclass[11pt]{article}
\usepackage{fontspec}
\usepackage[a4paper,margin=25mm]{geometry}
\usepackage{amsmath,amssymb,mathtools}
\usepackage{xurl}
\usepackage[hidelinks]{hyperref}
\newcommand{\sref}[2]{\hyperlink{source:#1:#2}{[S#2]}}
\setlength{\emergencystretch}{3em}
\begin{document}
% problemId:30007657
\section*{Aid that builds lasting capacity}\label{30007657}
\subsection{Definitions and mathematical statement}
Fix a set of aid-eligible communities and a program that ends at date \(t_0\). Let \(p\) specify funding, local participation, implementer, public-service responsibilities, and withdrawal rules. Let \(K_{it}(p)\) measure locally maintained administrative or productive capacity and \(Y_{it}(p)\) household economic welfare. For a declared horizon \(T>t_0\), define \(\tau_K(T,p)=E[K_{iT}(p)-K_{iT}(0)]\) and \(\tau_Y(T,p)=E[Y_{iT}(p)-Y_{iT}(0)]\). Capacity must be measured using verified service delivery, staffing, maintenance, or collective-action outputs, rather than continued external spending. Estimate whether aid complements state and community investment or substitutes for it, allowing political incentives and nonbeneficiary spillovers. Randomized introduction identifies offer effects; withdrawal comparisons need credible timing variation and controls for concurrent shocks. Separate temporary assets from institutional changes by measuring outcomes after support ends. Determine which program features predict durable gains and can be transported to other delivery systems. A known long-run null or positive result in one setting is a partial answer. The present statement proposes a general identification and design extension; no exact primary passage establishing this complete question as open was verified in the supporting long-run aid study.

\par\textbf{Formulation and scope.} This is a proposed scoped research specification. No primary passage explicitly stating this entire remaining question as open was verified. The supporting literature may answer important subquestions; this entry must not be advertised as a verified formal open problem.

\par\textbf{Open-question provenance.} An explicit published open-question statement was not verified. Sources below are supporting literature and must not be treated as a first listing.

\par\textbf{Historical priority.} No qualifying problem statement was directly verified, so historical priority cannot be assigned.

\subsection{Short English statement}
In the explicitly specified setting: Which aid designs produce durable economic gains and stronger local institutions after funding ends, and which create dependence or substitute for domestic capacity? The mathematical statement above fixes the target, assumptions and scope of this catalogue formulation.

\subsection{Sources}
\begin{itemize}\raggedright
\item[\textnormal{[S1]}]\hypertarget{source:30007657:1}{} Katherine Casey, Rachel Glennerster, Edward Miguel, Maarten J. Voors. 2021. Long Run Effects of Aid: Forecasts and Evidence from Sierra Leone. Abstract, paragraph describing eleven-year follow-up and institutional outcomes.
\url{https://www.nber.org/papers/w29079}.
DOI: 10.3386/w29079.
{\small\textit{Supporting or current literature; see evidence qualification. This follow-up supplies persistent infrastructure and market results but limited institutional change; the broad aid-versus-state-capacity question is not explicitly declared open.}}
\end{itemize}
\end{document}
